\documentclass[11pt,a4paper]{ctexart}
\usepackage[margin=23mm]{geometry}
\usepackage{amsmath,amssymb,fancyhdr}
\usepackage[colorlinks=true,linkcolor=blue,pdftitle={10.25作业2}]{hyperref}
\pagestyle{fancy}\fancyhf{}
\fancyhead[L]{微分方程}\fancyhead[R]{10.25作业2}
\fancyfoot[C]{\thepage}\setlength{\headheight}{15pt}
\setlength{\parindent}{0pt}\setlength{\parskip}{5pt}
\newcommand{\xx}{\boldsymbol{x}}
\newcommand{\bb}{\boldsymbol{b}}
\begin{document}
\begin{center}{\LARGE\bfseries 10.25作业2}\\[7pt]题目与解答\end{center}
\textbf{范围：}按用户提供的“10月21日作业”截图，解答第 110 页第 1 题 (1)、(3)、(4)，第 2 题 (1)、(3)，第 111 页第 7 题。题干及页码依据用户提供的截图。

\section*{第 110 页，第 1 题}
\textbf{题目。}求常系数齐次或非齐次线性系统 $\dfrac{dx}{dt}=Ax+b(t)$ 的通解，其中

\subsection*{(1)}
\[
A=\begin{pmatrix}0&-n^2\\-n^2&0\end{pmatrix},\qquad
b(t)=\begin{pmatrix}\cos nt\\\sin nt\end{pmatrix}.
\]
\textbf{解答。}先设实参数 $n\ne0$。矩阵 $A$ 的两个特征值及相应特征向量为
\[
\lambda_1=-n^2,\quad v_1=\begin{pmatrix}1\\1\end{pmatrix};\qquad
\lambda_2=n^2,\quad v_2=\begin{pmatrix}1\\-1\end{pmatrix}.
\]
故齐次系统的通解为
\[
\xx_h(t)=C_1e^{-n^2t}\begin{pmatrix}1\\1\end{pmatrix}
+C_2e^{n^2t}\begin{pmatrix}1\\-1\end{pmatrix}.
\]
对非齐次系统，试取特解
\[
\xx_p(t)=\begin{pmatrix}\alpha\sin nt\\\beta\cos nt\end{pmatrix}.
\]
代入方程得
\[
n\alpha=-n^2\beta+1,\qquad -n\beta=-n^2\alpha+1.
\]
解得
\[
\alpha=\frac{n+1}{n(n^2+1)},\qquad
\beta=\frac{n-1}{n(n^2+1)}.
\]
因此通解为
\[
\boxed{\xx(t)=C_1e^{-n^2t}\begin{pmatrix}1\\1\end{pmatrix}
+C_2e^{n^2t}\begin{pmatrix}1\\-1\end{pmatrix}
+\begin{pmatrix}
\dfrac{n+1}{n(n^2+1)}\sin nt\\[5pt]
\dfrac{n-1}{n(n^2+1)}\cos nt
\end{pmatrix}.}
\]
若允许 $n=0$，原系统变成 $x_1'=1,x_2'=0$，其通解为
\[
\boxed{\xx(t)=\begin{pmatrix}t+C_1\\C_2\end{pmatrix}.}
\]

\newpage
\subsection*{(3)}
\[
A=\begin{pmatrix}2&1&-2\\-1&0&0\\1&1&-1\end{pmatrix},\qquad
b(t)=\begin{pmatrix}0\\0\\0\end{pmatrix}.
\]
\textbf{解答。}特征多项式为
\[
\det(\lambda I-A)=(\lambda-1)(\lambda^2+1),
\]
故特征值为 $1,i,-i$。
对 $\lambda=1$ 可取特征向量
\[
v_1=\begin{pmatrix}1\\-1\\0\end{pmatrix}.
\]
对 $\lambda=i$ 可取特征向量
\[
v=\begin{pmatrix}1\\i\\1\end{pmatrix}
=\begin{pmatrix}1\\0\\1\end{pmatrix}
+i\begin{pmatrix}0\\1\\0\end{pmatrix}.
\]
因为 $A$ 为实矩阵，复解 $e^{it}v$ 的实部和虚部分别是实解：
\[
\operatorname{Re}(e^{it}v)=\begin{pmatrix}\cos t\\-\sin t\\\cos t\end{pmatrix},
\qquad
\operatorname{Im}(e^{it}v)=\begin{pmatrix}\sin t\\\cos t\\\sin t\end{pmatrix}.
\]
这两个解与 $e^tv_1$ 在 $t=0$ 时组成的矩阵为
\[
\begin{pmatrix}1&1&0\\-1&0&1\\0&1&0\end{pmatrix},
\qquad \det=-1\ne0,
\]
所以它们组成基本解组，通解为
\[
\boxed{\xx(t)=C_1e^t\begin{pmatrix}1\\-1\\0\end{pmatrix}
+C_2\begin{pmatrix}\cos t\\-\sin t\\\cos t\end{pmatrix}
+C_3\begin{pmatrix}\sin t\\\cos t\\\sin t\end{pmatrix}.}
\]

\subsection*{(4)}
\[
A=\begin{pmatrix}0&1&0&0\\-1&0&0&0\\0&0&0&-1\\0&0&1&0\end{pmatrix},\qquad
b(t)=\begin{pmatrix}0\\0\\0\\0\end{pmatrix}.
\]
\textbf{解答。}系统分成两组独立的方程
\[
x_1'=x_2,\quad x_2'=-x_1;\qquad
x_3'=-x_4,\quad x_4'=x_3.
\]
由 $x_1''+x_1=0$、$x_3''+x_3=0$ 及上述一阶关系，得到通解
\[
\boxed{\xx(t)=
\begin{pmatrix}
\cos t&\sin t&0&0\\
-\sin t&\cos t&0&0\\
0&0&\cos t&-\sin t\\
0&0&\sin t&\cos t
\end{pmatrix}
\begin{pmatrix}C_1\\C_2\\C_3\\C_4\end{pmatrix}.}
\]
所列矩阵在 $t=0$ 时为 $I_4$，是该系统的标准基本解矩阵。

\newpage
\section*{第 110 页，第 2 题}
\textbf{题目。}求微分系统 $\dfrac{dx}{dt}=Ax+b(t)$ 满足初值条件 $x(0)=\eta$ 的解，其中

\subsection*{(1)}
\[
A=\begin{pmatrix}-5&-1\\1&-3\end{pmatrix},\qquad
b(t)=\begin{pmatrix}e^t\\e^{2t}\end{pmatrix},\qquad
\eta=\begin{pmatrix}1\\0\end{pmatrix}.
\]
\textbf{解答。}写成
\[
A=-4I+N,\qquad N=\begin{pmatrix}-1&-1\\1&1\end{pmatrix},\qquad N^2=0.
\]
因此，由矩阵指数的幂级数，$e^{Nt}=I+tN$，故
\[
e^{At}=e^{-4t}(I+tN)
=e^{-4t}\begin{pmatrix}1-t&-t\\t&1+t\end{pmatrix}.
\]
试取非齐次特解 $\xx_p(t)=pe^t+qe^{2t}$。代入原系统并比较指数项，得到
\[
(I-A)p=\begin{pmatrix}1\\0\end{pmatrix},\qquad
(2I-A)q=\begin{pmatrix}0\\1\end{pmatrix}.
\]
即
\[
\begin{pmatrix}6&1\\-1&4\end{pmatrix}p=\begin{pmatrix}1\\0\end{pmatrix},\qquad
\begin{pmatrix}7&1\\-1&5\end{pmatrix}q=\begin{pmatrix}0\\1\end{pmatrix}.
\]
解得
\[
p=\frac1{25}\begin{pmatrix}4\\1\end{pmatrix},\qquad
q=\frac1{36}\begin{pmatrix}-1\\7\end{pmatrix}.
\]
因此，满足初值的解可写成
\[
\xx(t)=e^{At}c+\frac{e^t}{25}\begin{pmatrix}4\\1\end{pmatrix}
+\frac{e^{2t}}{36}\begin{pmatrix}-1\\7\end{pmatrix},
\]
其中
\[
c=\eta-p-q=\frac1{900}\begin{pmatrix}781\\-211\end{pmatrix},\qquad
Nc=\frac1{900}\begin{pmatrix}-570\\570\end{pmatrix}.
\]
所以最终答案为
\[
\boxed{\xx(t)=\frac{e^{-4t}}{900}
\begin{pmatrix}781-570t\\-211+570t\end{pmatrix}
+\frac{e^t}{25}\begin{pmatrix}4\\1\end{pmatrix}
+\frac{e^{2t}}{36}\begin{pmatrix}-1\\7\end{pmatrix}.}
\]

\newpage
\subsection*{(3)}
\[
A=\begin{pmatrix}4&-3\\2&-1\end{pmatrix},\qquad
b(t)=\begin{pmatrix}\sin t\\-2\cos t\end{pmatrix},\qquad
\eta=\begin{pmatrix}0\\0\end{pmatrix}.
\]
\textbf{解答。}由
\[
\det(\lambda I-A)=(\lambda-1)(\lambda-2),
\]
得两个特征值 $1,2$，对应特征向量分别可取 $(1,1)^{\mathsf T}$、$(3,2)^{\mathsf T}$。
因此齐次通解为
\[
\xx_h(t)=C_1e^t\begin{pmatrix}1\\1\end{pmatrix}
+C_2e^{2t}\begin{pmatrix}3\\2\end{pmatrix}.
\]
试取特解
\[
\xx_p(t)=u\sin t+v\cos t.
\]
代入并比较 $\sin t$、$\cos t$ 的系数，得到
\[
-v=Au+\begin{pmatrix}1\\0\end{pmatrix},\qquad
u=Av+\begin{pmatrix}0\\-2\end{pmatrix}.
\]
消去 $v$，有
\[
(I+A^2)u=-A\begin{pmatrix}1\\0\end{pmatrix}
+\begin{pmatrix}0\\-2\end{pmatrix}
=\begin{pmatrix}-4\\-4\end{pmatrix}.
\]
由于
\[
I+A^2=\begin{pmatrix}11&-9\\6&-4\end{pmatrix},
\]
解得
\[
u=\begin{pmatrix}-2\\-2\end{pmatrix},\qquad
v=-Au-\begin{pmatrix}1\\0\end{pmatrix}=\begin{pmatrix}1\\2\end{pmatrix}.
\]
于是通解为
\[
\xx(t)=C_1e^t\begin{pmatrix}1\\1\end{pmatrix}
+C_2e^{2t}\begin{pmatrix}3\\2\end{pmatrix}
-2\sin t\begin{pmatrix}1\\1\end{pmatrix}
+\cos t\begin{pmatrix}1\\2\end{pmatrix}.
\]
利用 $\xx(0)=0$，得
\[
C_1+3C_2+1=0,\qquad C_1+2C_2+2=0,
\]
故 $C_1=-4,C_2=1$。所求解为
\[
\boxed{\xx(t)=-4e^t\begin{pmatrix}1\\1\end{pmatrix}
+e^{2t}\begin{pmatrix}3\\2\end{pmatrix}
-2\sin t\begin{pmatrix}1\\1\end{pmatrix}
+\cos t\begin{pmatrix}1\\2\end{pmatrix}.}
\]

\newpage
\section*{第 111 页，第 7 题}
\textbf{题目。}求出所有 2 阶常数矩阵的实 Jordan 标准形及其相应的齐次线性系统的基本解矩阵。

\textbf{解答。}按实矩阵理解题中的“实 Jordan 标准形”。任意 $A\in\mathbb R^{2\times2}$ 均可用实可逆矩阵 $P$ 化成 $A=PJP^{-1}$。令 $\xx=P\boldsymbol{y}$，则
\[
\xx'=A\xx\quad\Longleftrightarrow\quad \boldsymbol{y}'=J\boldsymbol{y}.
\]
以下逐一列出所有可能的 $J$，以及满足 $\Psi(0)=I$ 的标准基本解矩阵 $\Psi(t)=e^{Jt}$。

\subsection*{一、两个不同的实特征值}
\[
\boxed{J=\begin{pmatrix}\lambda_1&0\\0&\lambda_2\end{pmatrix},\qquad
\Psi(t)=\begin{pmatrix}e^{\lambda_1t}&0\\0&e^{\lambda_2t}\end{pmatrix},
\quad \lambda_1,\lambda_2\in\mathbb R,\ \lambda_1\ne\lambda_2.}
\]
可取两个实特征向量为 $P$ 的列向量；交换特征值的次序给出同一类型。

\subsection*{二、一个二重实特征值，且有两个线性无关的特征向量}
\[
\boxed{J=\begin{pmatrix}\lambda&0\\0&\lambda\end{pmatrix}=\lambda I,
\qquad \Psi(t)=e^{\lambda t}I,\quad \lambda\in\mathbb R.}
\]
此时 $A$ 本身就是 $\lambda I$。

\subsection*{三、一个二重实特征值，且只有一个线性无关的特征向量}
\[
\boxed{J=\begin{pmatrix}\lambda&1\\0&\lambda\end{pmatrix},\qquad
\Psi(t)=e^{\lambda t}\begin{pmatrix}1&t\\0&1\end{pmatrix},\quad\lambda\in\mathbb R.}
\]
这里 $J=\lambda I+N$，$N^2=0$，故 $e^{Jt}=e^{\lambda t}(I+tN)$。
可取特征向量 $p$ 和广义特征向量 $q$，满足
\[
(A-\lambda I)p=0,\qquad (A-\lambda I)q=p,
\]
并令 $P=(p,q)$。

\subsection*{四、一对共轭非实特征值}
设特征值为 $a\pm bi$，$a\in\mathbb R,b>0$，则可取
\[
\boxed{J=\begin{pmatrix}a&-b\\b&a\end{pmatrix},\qquad
\Psi(t)=e^{at}\begin{pmatrix}\cos bt&-\sin bt\\\sin bt&\cos bt\end{pmatrix}.}
\]
事实上，写成 $J=aI+bK$，其中
\[
K=\begin{pmatrix}0&-1\\1&0\end{pmatrix},\qquad K^2=-I,
\]
由幂级数得 $e^{bKt}=I\cos bt+K\sin bt$，即得上述公式。
若 $p+iq$ 是 $a+bi$ 的特征向量，则可取 $P=(p,-q)$。
另一常用约定是取 $\widetilde J=\begin{pmatrix}a&b\\-b&a\end{pmatrix}$，它与上述 $J$ 实相似，此时对应的基本解矩阵为
\[
e^{at}\begin{pmatrix}\cos bt&\sin bt\\-\sin bt&\cos bt\end{pmatrix}.
\]

\textbf{回到原系统。}对以上任一类型，$\Psi'=J\Psi$ 且 $\Psi(0)=I$，所以
\[
\boxed{\Phi(t)=P\Psi(t)}
\]
是原系统 $\xx'=A\xx$ 的基本解矩阵；若要求在 $t=0$ 处等于 $I$，则取
\[
\boxed{X(t)=P\Psi(t)P^{-1}=e^{At}.}
\]
实二次特征多项式的根只可能是两个不同实根、一个二重实根或一对共轭非实根；二重实根再按特征空间维数分为上述两种情形，因此分类穷尽。
\end{document}
